\subsection{The exact maximal solvable quotient of the first Hurwitz pair}

\begin{statement}{D23.F35}{Proved}{Perfect residual and a marked solvable quotient of order $93312$}
\end{statement}
In the notation of D23.F34, the full pure-pair image
$P=\langle A^2,B^2\rangle$ has the following exact derived series:
\[
\begin{array}{c|r}
 i&|P^{(i)}|\\\hline
0&1567283281920000\\
1&261213880320000\\
2&16325867520000\\
3&4081466880000\\
4&50388480000\\
5&16796160000\\
6&16796160000.
\end{array}                                               \tag{D23.F35.1}
\]
In particular its perfect residual is
\[
 N=P^{(5)}=P^{(6)},\qquad |N|=16796160000=360^4.
\]
This order identity is not an assertion identifying $N$ with a
product of alternating groups.

Restrict the first pair to the $448$-point subset of D23.P21, and
simultaneously take its signed degree-$24$ quotient from D23.F30.
Their marked joint permutation image $Q$ has degree $472$, order
$93312$, and derived series
\[
               93312,15552,972,243,3,1.                  \tag{D23.F35.2}
\]
The natural marked surjection $P\twoheadrightarrow Q$ has kernel
exactly $N$. Therefore
\[
                Q\cong P/P^{(\infty)},\qquad
                |P'/P^{(\infty)}|=15552.
\]
Thus this concrete joint image is the maximal solvable quotient of
the full first-pair image, with its original two free generators
retained.

\begin{proof}
Start with the literal degree-$1344$ generators of D23.P17 and form
$P$ from their squares. At each stage, form the normal subgroup
generated by all commutators of the current generators and compute
its exact order with a stabilizer chain. This gives precisely the
orders in (D23.F35.1). When a generator list is unnecessarily long,
the certificate replaces it by a shorter list of explicit products
in those generators. The replacement subgroup is contained in the
previous one by construction, and equality of their computed finite
orders proves equality of the groups. The random choice of products
is only an implementation aid: no probabilistic conclusion about
generation or group recognition is used.

The sixth derived subgroup is contained in the fifth. Their equal
finite orders prove equality, so the fifth is perfect and the series
has stabilized. This proves the assertion about $N$, including that
it is the intersection of the whole derived series.

For each pure generator, concatenate its restriction to the ordered
$448$-point invariant subset with its signed action on $24$ points.
Both are actual quotients of $P$: the former is restriction to an
invariant subset, and the latter is the signed quotient proved in
D23.F30. Hence the resulting $472$-point group is a marked quotient
$Q$ of $P$, without any presentation or order-based identification.
The literal two arrays are checked against the separately
constructed root record
\path{work/records/hurwitz_first_solvable_joint.npz}.
Computing all its successive derived groups gives exactly
(D23.F35.2), and in particular proves that $Q$ is solvable.

A homomorphism sends every term of a derived series into the
corresponding term of the target series. Since $Q^{(5)}=1$, the
kernel of $P\to Q$ contains $P^{(5)}=N$. On the other hand,
\[
        |\ker(P\to Q)|=\frac{1567283281920000}{93312}
                      =16796160000=|N|.
\]
Thus the kernel equals $N$. Finally the image of the commutator
subgroup under a surjection is the target commutator subgroup,
so $P'/N\cong Q'$ has order $15552$. Every solvable quotient of
$P$ kills its perfect residual; consequently $Q=P/N$ has the stated
maximal property.
\end{proof}

\paragraph{Exact replay and scope.}
The complete certificate is
\path{agents/five_level_six/hurwitz_first_perfect_residual.py},
with JSON and NPZ records of the same stem. It starts from the
original $1344$-point generators, constructs every derived group,
checks stabilization, constructs the marked $472$-point quotient,
and checks its complete derived series. No cached group order is
accepted as an input. The NPZ retains generating lists at every
step and the literal two quotient generators. This statement
identifies the maximal solvable quotient only; common-word lifting
and actual gentle or full-GT realization are separate assertions.
